\article{论全纯尖点形式的提升}{On liftings of holomorphic cusp forms}{Takuro Shintani}{2}{97--110}

\section*{引言（Introduction）}
在 [10] 中，H. Saito 计算了“扭曲”Hecke 算子（twisted Hecke operator）的迹，这些算子作用于素数次循环全实阿贝尔域上的 Hilbert 模群（Hilbert modular group）的全纯尖点形式（holomorphic cusp form）空间。他发现，这一“扭曲迹公式”（twisted trace formula）与椭圆模形式（elliptic modular form）空间上通常的 Hecke 迹公式之间，存在一个引人注目的恒等关系。利用此恒等式，他证明了椭圆模形式可以“提升”（lifting）为 Hilbert 模形式（Hilbert modular form）（关于此类提升结果的起源，见 [2]）。他还进一步刻画了所有提升形式张成的空间。在 1975 年 6 月于 Ann Arbor 举行的美日研讨会“自守形式在数论中的应用”上，作者报告了 Saito 工作的一种表示论解释及推广（见 [13]）。Langlands 随即在 [8] 中推广了作者报告中的结果。此外，Langlands 还发现了“提升理论”在 Artin $L$ 函数理论中的一个出人意料的应用。

本文分为两节：第 1 节在稍作调整后重述 Ann Arbor 报告的内容，第 2 节补充若干讨论。

作者感谢 H. Saito 教授在 [10] 出版之前详细讲解其中的工作。作者此前的两篇文章 [14]、[15] 也受到 [10] 的启发。\footnote{原文的 AMS（MOS）主题分类（1970）：主要分类 10Dxx；次要分类 10D15、10D20。原文版权：1979，American Mathematical Society。}

\paragraph*{记号（Notation）.}
对环 $R$，$R^\times$ 表示其单位群（group of units）。记 $G(R)=G_R=\GL(2,R)$。对 $p$ 进域（$p$-adic field）$k$，$\mathfrak o(k)=\mathfrak o_k$ 表示 $k$ 的整数环（ring of integers），$q_k$ 表示其剩余域（residue class field）的基数。对每个 $t\in k$，令 $d(tx)=\alpha_k(t)\,dx=|t|_k\,dx$，其中 $dx$ 是 $k$ 的加法群的不变测度（invariant measure）。对每个局部紧全不连通群（locally compact totally disconnected group）$G$，$C_0^\infty(G)$ 是 $G$ 上的紧支集局部常值函数空间。设 $k$ 是域，$F$ 是 $k$ 的扩域，用 $N_k^F$ 表示从 $F$ 到 $k$ 的范数映射（norm mapping）。

\origsection{1}{}
\origsubsection{1.1}{}
设 $k$ 为有限代数数域或 $p$ 进域。设 $F$ 为 $k$ 上素数次 $l$ 的交换半单代数（commutative semisimple algebra），并假定 $F$ 在 $k$ 上的自同构群包含一个由 $\sigma$ 生成的 $l$ 阶循环子群 $\mathfrak g$。那么 $F$ 或者同构于 $k$ 的一个循环扩域，或者同构于 $l$ 个 $k$ 的直和。令 $G_F=\GL(2,F)$，$G_k=\GL(2,k)$。可自然地将 $\mathfrak g$ 视为 $G_F$ 的一个自同构群，其不动点集为 $G_k$。若存在 $g\in G_F$ 使 $y=g^\sigma xg^{-1}$，便称 $G_F$ 中的元素 $x,y$ 在 $G_F$ 中为 $\sigma$-扭曲共轭（$\sigma$-twistedly conjugate）。用 $x^{G_F,\sigma}$ 表示 $G_F$ 中所有与 $x$ 为 $\sigma$-扭曲共轭的元素的集合，用 $\mathcal C_\sigma(G_F)$ 表示 $G_F$ 的 $\sigma$-扭曲共轭类集合。包含 $x$ 的通常共轭类记为 $x^{G_F}$。对每个 $x\in G_F$，令
\[
 N(x)=x^{\sigma^{l-1}}x^{\sigma^{l-2}}\cdots x^\sigma x.
\]
容易看出 $N(g^\sigma xg^{-1})=gN(x)g^{-1}$。

用 $\mathcal C(G_k)$ 表示 $G_k$ 的共轭类集合。下述引理的证明见 [10] 的引理 3.4 与引理 3.6。

\begin{originaltheorem}{引理 1（Lemma 1，Saito）}
沿用上述记号与假设，则映射
\[
 x^{G_F,\sigma}\longmapsto N(x)^{G_F}\cap G_k
\]
是从 $\mathcal C_\sigma(G_F)$ 到 $\mathcal C(G_k)$ 的单射。若 $F$ 不是域，则此映射是双射。
\end{originaltheorem}

下文中，对每个 $c\in\mathcal C_\sigma(G_F)$，用 $N(c)\in\mathcal C(G_k)$ 表示它在引理 1 的映射下的像。如果 $x\in G_k$ 在 $G_k$ 中的中心化子（centralizer）是一个二维 $k$-环面（$k$-torus），则称 $x$ 为正则的（regular）。若 $G_k$ 的某个共轭类由正则元素组成，则称此共轭类正则。若 $N(c)$ 正则，则称 $G_F$ 的 $\sigma$-扭曲共轭类 $c$ 正则。分别用 $\mathcal C'_\sigma(G_F)$ 与 $\mathcal C'(G_k)$ 表示 $G_F$ 的正则 $\sigma$-扭曲共轭类集合与 $G_k$ 的正则共轭类集合。

\origsubsection{1.2}{}
沿用 1.1 的记号，并假定 $k$ 为 $p$ 进域。令 $\widetilde G_F$ 为 $G_F$ 与 $\mathfrak g$ 的半直积（semidirect product）。确切地说，其底集为 $\mathfrak g\times G_F$，乘法法则为
\[
 (\tau,g)(\tau',g')=(\tau\tau',g^{\tau'}g'),
 \qquad\tau,\tau'\in\mathfrak g,\quad g,g'\in G_F.
\]
由引理 1 立即可知，$(\sigma,g_1)$ 与 $(\sigma,g_2)$ 在 $\widetilde G_F$ 中共轭，当且仅当 $g_1$ 与 $g_2$ 在 $G_F$ 中为 $\sigma$-扭曲共轭。用 $\widehat G_k$（分别为 $\widehat G_F$）表示 $G_k$（分别为 $G_F$）的不可约容许可酉化表示（irreducible admissible unitarizable representation）的等价类集合。对每个 $R\in\widehat G_F$ 和 $\tau\in\mathfrak g$，令 $R^\tau$ 为由 $R^\tau(x)=R(x^\tau)$ 给出的 $G_F$ 的表示。于是 $R^\tau\in\widehat G_F$，从而 $\mathfrak g$ 作用于 $\widehat G_F$。用 $\widehat G_F^{\mathfrak g}$ 表示其 $\mathfrak g$-不动元素子集。若 $\widetilde G_F$ 的一个表示在 $G_F$ 上的限制是容许的，便称它容许。设 $\widetilde R$ 是 $\widetilde G_F$ 的不可约容许可酉化表示。众所周知，它在 $G_F$ 上的限制 $R$，或者不可约，或者是 $l$ 个互不等价的不可约表示的直和。在前一种情形下，称 $\widetilde R$ 为第一类（of the first kind）。假定 $\widetilde R$ 为第一类，令 $J_\sigma=\widetilde R((\sigma,1))$，则
\begin{equation}
 R(g^\sigma)=J_\sigma^{-1}R(g)J_\sigma,\qquad g\in G_F.
 \tag{1.1}
\end{equation}
因此 $R\in\widehat G_F^{\mathfrak g}$。关系 (1.1) 在相差一个 $l$ 次单位根倍数的范围内确定了算子 $J_\sigma$。反过来，若 $R\in\widehat G_F^{\mathfrak g}$，则存在满足 (1.1) 的 $l$ 阶酉算子。于是令 $\widetilde R((\sigma^m,g))=J_\sigma^mR(g)$，便将 $R$ 延拓为 $\widetilde G_F$ 的第一类容许不可约可酉化表示。对 $\varphi\in C_0^\infty(G_F)$，用 $\widetilde R(\varphi)$ 表示算子
\[
 \widetilde R(\varphi)=\int_{G_F}\widetilde R((\sigma,g))\varphi(g)\,dg,
\]
其中 $dg$ 是 $G_F$ 上的不变测度，正规化为 $G_{\mathfrak o(F)}$ 的体积等于 $1$。

可以证明，$G_F$ 上存在局部可积函数（locally summable function）$\chi(\widetilde R)$，使
\[
 \operatorname{trace}\widetilde R(\varphi)
 =\int_{G_F}\varphi(g)\chi(\widetilde R)(g)\,dg.
\]
还可证明，$\chi(\widetilde R)(g)$ 仅依赖于 $g$ 在 $G_F$ 中的 $\sigma$-扭曲共轭类，且 $\chi(\widetilde R)$ 只在 $\mathcal C'_\sigma(G_F)$ 上定义。对每个 $r\in\widehat G_k$，用 $\chi_r$ 表示其特征标（character）。回忆 $\chi_r(x)$ 仅依赖于 $x$ 的共轭类，且 $\chi_r$ 只在 $\mathcal C'(G_k)$ 上定义。进一步假定 $k$ 的剩余域特征不等于 $2$。

\begin{originaltheorem}{定理 1（Theorem 1）}
沿用上述记号与假设。

(1) 对每个 $R\in\widehat G_F^{\mathfrak g}$，可将 $R$ 延拓为 $\widetilde G_F$ 的表示 $\widetilde R$，并可找到 $r\in\widehat G_k$，使
\begin{equation}
 \chi(\widetilde R)(x)=\chi_r(N(x)),
 \tag{1.2}
\end{equation}
其中 $N$ 为引理 1 给出的从 $\mathcal C'_\sigma(G_F)$ 到 $\mathcal C'(G_k)$ 的范数映射。

(2) 对每个 $r\in\widehat G_k$，存在唯一的 $R\in\widehat G_F^{\mathfrak g}$，使得 $R$ 适当延拓到 $\widetilde G_F$ 后，所得 $\widetilde R$ 满足 (1.2)。
\end{originaltheorem}

对每个 $r\in\widehat G_k$，将由 (1.2) 与 $r$ 相关联的 $R\in\widehat G_F^{\mathfrak g}$ 称为 $r$ 从 $\widehat G_k$ 到 $\widehat G_F^{\mathfrak g}$ 的提升（lifting）。（Jacquet 在 [4] 中从不同观点引入了提升的概念。）

\paragraph*{注记 1.}
[14] 给出了有限一般线性群（finite general linear group）的定理 1 的类似结果。此外，[15]（分别为 [8]）用局部（分别为整体）方法给出了 $(F,k)=(\C,\R)$ 情形下定理 1 的类似结果。

下面具体说明 $\widehat G_k$ 到 $\widehat G_F^{\mathfrak g}$ 的提升如何构造。若 $F$ 作为 $k$-代数同构于 $l$ 个 $k$ 的直和，则 $G_F$ 同构于 $l$ 个 $G_k$ 的直积。容易看出，$r\in\widehat G_k$ 的提升为 $r\otimes\cdots\otimes r$（即使 $k$ 的剩余域特征等于 $2$，这也成立）。

接着考虑 $F$ 是 $k$ 的素数次 $l$ 的循环扩域的情形。首先回顾 $\widehat G_k$ 的描述，详见 [5] 第 1 章第 3、4 节。对满足 $\mu_1\mu_2^{-1}\ne\alpha_k^{\pm1}$ 的 $k^\times$ 的拟特征（quasi-character）$\mu_1,\mu_2$，令 $\rho(\mu_1,\mu_2)$ 为对应的 $G_k$ 的不可约主系列表示（principal series representation）。若 $\mu_1\mu_2^{-1}=\alpha_k$，令 $\sigma(\mu_1,\mu_2)$ 为对应的特殊表示（special representation）。对于 $k$ 的二次扩域 $K$ 和满足 $\omega'\ne\omega$ 的 $K^\times$ 的拟特征 $\omega$（$'$ 表示 $K$ 关于 $k$ 的共轭），令 $\pi(\omega,K)$ 为对应的绝对尖点表示（absolutely cuspidal representation）。若 $k$ 的剩余域特征不等于 $2$（正如当前假设），则已知 $G_k$ 的每个无限维不可约容许表示都等价于某个 $\rho(\mu_1,\mu_2)$、$\sigma(\mu_1,\mu_2)$ 或 $\pi(\omega,K)$。对 $k^\times$ 的每个拟特征 $\mu$，用 $\boldsymbol\mu$ 表示 $F^\times$ 的拟特征 $\boldsymbol\mu=\mu\circ N_k^F$。对于 $k$ 的二次扩域 $K\ne F$ 和 $K^\times$ 的拟特征 $\omega$，用 $\boldsymbol\omega$ 表示 $L^\times$（$L=K\cdot F$）的拟特征 $\boldsymbol\omega=\omega\circ N_K^L$。

\begin{originaltheorem}{命题 1（Proposition 1）}
沿用上述记号与假设，$r\in\widehat G_k$ 到 $\widehat G_F^{\mathfrak g}$ 的提升 $R$ 如下：

(1) 若 $r=\rho(\mu_1,\mu_2)$（分别为 $\sigma(\mu_1,\mu_2)$），则 $R=\rho(\boldsymbol\mu_1,\boldsymbol\mu_2)$（分别为 $R=\sigma(\boldsymbol\mu_1,\boldsymbol\mu_2)$）。

(2) 若 $r=\pi(\omega,K)$ 且 $K\ne F$，则 $R=\pi(\boldsymbol\omega,L)$，其中 $L=F\cdot K$。\footnote{\textbf{译者注（本次校订）：}按前文只对 $\omega\ne\omega^{\prime}$ 定义尖点表示 $\pi(\omega,K)$ 的约定，这里还须检查 $\boldsymbol\omega\ne\boldsymbol\omega^{\prime}$。当 $l=2$ 时，$\omega/\omega^{\prime}$ 可能在 $N_K^L(L^\times)$ 上平凡，从而使 $\boldsymbol\omega$ 在 $L/F$ 的共轭下不变。此时相应的二维诱导参数分裂，提升是由 $\boldsymbol\omega$ 的两个延拓所确定的主系列，而非尖点表示。}

(3) 若 $r=\pi(\omega,K)$ 且 $K=F$，则 $R=\rho(\omega,\omega')$。
\end{originaltheorem}

\paragraph*{注记 2.}
命题 1 的第一项可直接验证（即使 $k$ 的剩余域特征等于 $2$ 也是如此）。

\origsubsection{1.3}{}
沿用 1.2 的记号。特别地，$k$ 为 $p$ 进域（其剩余域可以具有特征 $2$）。对 $x\in G_F$，令 $Z_\sigma(x)$ 为 $G_F$ 中满足 $g^\sigma xg^{-1}=x$ 的所有元素 $g$ 构成的子群。将 $G_F$ 和 $Z_\sigma(x)$ 上的不变测度正规化，使各自的极大紧开子群（maximal compact open subgroup）的总体积均等于 $1$。用 $d\dot g$ 表示由这两个正规化测度的商给出的 $G_F/Z_\sigma(x)$ 上的不变测度。对 $f\in C_0^\infty(G_F)$，令
\begin{equation}
 \Lambda_\sigma(f,x)=\int_{G_F/Z_\sigma(x)}f(g^\sigma xg^{-1})\,d\dot g.
 \tag{1.3}
\end{equation}
可以证明此积分绝对收敛。此外，对 $c\in\mathcal C'_\sigma(G_F)$，任意 $x,y\in c$ 均满足 $\Lambda_\sigma(f,x)=\Lambda_\sigma(f,y)$。于是对每个 $c\in\mathcal C'_\sigma(G_F)$，令
\begin{equation}
 \Lambda_\sigma(f,c)=\Lambda_\sigma(f,x),
 \tag{1.4}
\end{equation}
其中 $x$ 是 $c$ 的任意元素。对每个 $x\in G_k$，令 $Z(x)$ 为 $x$ 在 $G_k$ 中的中心化子。将 $G_k$ 和 $Z(x)$ 上的不变测度正规化，使各自的极大紧开子群的总体积均等于 $1$。用 $d\dot g$ 表示由这两个正规化测度的商给出的 $G_k/Z(x)$ 上的不变测度。对 $f\in C_0^\infty(G_k)$，令
\begin{equation}
 \Lambda(f,x)=\int_{G_k/Z(x)}f(gxg^{-1})\,d\dot g.
 \tag{1.5}
\end{equation}
已知此积分绝对收敛。对每个 $c\in\mathcal C'(G_k)$，令
\begin{equation}
 \Lambda(f,c)=\Lambda(f,x),
 \tag{1.6}
\end{equation}
其中 $x$ 是 $c$ 的任意元素。容易看出 (1.6) 的右端与 $x\in c$ 的选择无关。

\begin{originaltheorem}{命题 2（Proposition 2）}
沿用上述记号，则对每个 $\boldsymbol f\in C_0^\infty(G_F)$，存在 $f\in C_0^\infty(G_k)$ 满足：

(1) 对每个正则的 $c\in\mathcal C_\sigma(G_F)$，$\Lambda_\sigma(\boldsymbol f,c)=\Lambda(f,N(c))$。

(2) 对每个 $c\in\mathcal C'(G_k)\setminus N(\mathcal C'_\sigma(G_F))$，$\Lambda(f,c)=0$。
\end{originaltheorem}

\paragraph*{注记 3.}
在命题 2 中，假设 $\boldsymbol f$ 为 $G_{\mathfrak o(F)}$ 的示性函数（characteristic function），并进一步假设 $F$ 或者为 $k$ 的 $l$ 次非分歧循环扩域，或者同构于 $l$ 个 $F$ 的直和，则可以取 $f$ 为 $G_{\mathfrak o(k)}$ 的示性函数。

\noindent\emph{译注：所供扫描本注记 3 末句为“$l$ 个 $F$ 的直和”；按 1.1 的分裂情形，通常应为“$l$ 个 $k$ 的直和”。这里保留扫描本的域名。}

\paragraph*{注记 4.}
命题 2 的函数对应 $\boldsymbol f\mapsto f$，与定理 1 的不可约特征标提升有一种对偶关系。

\origsubsection{1.4}{}
设 $k$ 为 $n$ 次全实代数数域，用 $k_\A$（分别为 $k_\A^\times$）表示 $k$ 的阿代尔环（adele ring）（分别为伊代尔群（idele group））。用 $k_\infty$（分别为 $k_{\A,0}$）表示 $k_\A$ 的无穷分量（分别为有限分量）。于是 $k_\A=k_\infty\oplus k_{\A,0}$，$k_\A^\times=k_\infty^\times\times k_{\A,0}^\times$。此外，$G_{k_\infty}$ 和 $G_{k_{\A,0}}$ 均自然地嵌入 $G_{k_\A}$，且 $G_{k_\A}=G_{k_\infty}\times G_{k_{\A,0}}$。设 $\{\infty_1,\ldots,\infty_n\}$ 为 $k$ 的无穷素位集合。对每个 $g\in G_{k_\A}$，写 $g=g_\infty g_0$，其中 $g_\infty\in G_{k_\infty}$，$g_0\in G_{k_{\A,0}}$，并写 $g_\infty=(g_{\infty,1},\ldots,g_{\infty,n})$，其中 $g_{\infty,i}\in G_\R$ 为对应于无穷素位 $\infty_i$ 的分量。对 $g_{\infty,i}$ 引入标准参数化
\[
 g_{\infty,i}=
 \begin{pmatrix}t_i&0\\0&\pm t_i\end{pmatrix}
 \begin{pmatrix}1&x_i\\0&1\end{pmatrix}
 \begin{pmatrix}\sqrt{y_i}&0\\0&\sqrt{y_i^{-1}}\end{pmatrix}
 \begin{pmatrix}\cos\theta_i&\sin\theta_i\\-\sin\theta_i&\cos\theta_i\end{pmatrix},
\]
其中 $t_i>0$，$x_i\in\R$，$y_i>0$，$\theta_i\in\R$。

令 $\Omega_i$ 为 $G_{k_\A}$ 上由下式给出的微分算子（differential operator）：
\[
 \Omega_i=-y_i\frac{\partial^2}{\partial\theta_i\partial x_i}
 +y_i^2\left(\frac{\partial^2}{\partial x_i^2}
 +\frac{\partial^2}{\partial y_i^2}\right).
\]
对 $\theta=(\theta_1,\ldots,\theta_n)\in\R^n$，用 $\kappa(\theta)$ 表示 $G_{k_\infty}$ 中第 $i$ 个分量为
\[
 \begin{pmatrix}\cos\theta_i&\sin\theta_i\\-\sin\theta_i&\cos\theta_i\end{pmatrix}
\]
的元素。设 $\chi$ 为群 $k_\A^\times/k^\times$ 的酉特征（unitary character），$r$ 为 $k$ 的无穷素位集合上取正整数值的函数，令 $r_i=r(\infty_i)$。用 $S(\chi,r,k_\A)$ 表示 $G_{k_\A}$ 上满足以下条件的复值函数 $f$ 的空间：
\begin{equation}
 \begin{gathered}
 f(g)\text{ 有界，关于 }g_\infty\text{ 光滑，关于 }g_0\text{ 局部常值；}\\
 \Omega_i f=f\,r_i(r_i-2)/4,\qquad 1\le i\le n.
 \end{gathered}
 \tag{1.7}
\end{equation}
\begin{equation}
 f\left(\gamma\begin{pmatrix}z&0\\0&z\end{pmatrix}g\right)
 =\chi(z)f(g),\qquad\gamma\in G_k,\quad z\in k_\A^\times.
 \tag{1.8}
\end{equation}
\begin{equation}
 f(g\kappa(\theta))=\exp\left(\sqrt{-1}\sum_{i=1}^n r_i\theta_i\right)f(g),
 \qquad\theta\in\R^n.
 \tag{1.9}
\end{equation}

用 $\mathfrak p(k)$ 表示 $k$ 的有限素理想集合。对每个 $v\in\mathfrak p(k)$，令 $k_v$ 为 $k$ 在 $v$ 处的完备化，$\mathfrak o_v$ 为 $k_v$ 的整数环。于是 $G_{k_{\A,0}}$ 同构于限制直积（restricted direct product）$\prod_v'G_{k_v}$。将其 Haar 测度正规化，使 $\prod_vG_{\mathfrak o_v}$ 的体积为 $1$。对每个 $\varphi\in C_0^\infty(G_{k_{\A,0}})$，用 $T(\varphi)$ 表示 $S(\chi,r,k_\A)$ 上的线性算子
\[
 (T(\varphi)f)(g)=\int_{G_{k_{\A,0}}}f(gx)\varphi(x)\,dx.
\]
可证明 $T(\varphi)$ 具有有限秩。设 $F$ 是 $k$ 的素数次 $l$ 的全实循环扩域。令 $\boldsymbol\chi=\chi\circ N_{k_\A}^{F_\A}$，则 $\boldsymbol\chi$ 为 $F_\A^\times/F^\times$ 的特征。令 $r(w)=r(\overline w)$，将 $r$ 延拓为 $F$ 的无穷素位集合上的函数，其中 $w$ 是 $F$ 的任意无穷素位，$\overline w$ 为其在 $k$ 上的限制。

对每个 $\tau\in\mathfrak g=\Gal(F/k)$ 及 $G_{F_\A}$ 上的任意函数 $f$，用 $J_\tau f$ 表示由 $(J_\tau f)(g)=f(g^\tau)$ 定义的函数。容易看出，$J_\tau$ 保持空间 $S(\boldsymbol\chi,r,F_\A)$ 不变。对每个 $v\in\mathfrak p(k)$，令 $F_v=k_v\otimes_kF$，$\mathfrak o(F_v)$ 为 $F_v$ 的极大紧子环。群 $\mathfrak g$ 自然作用于 $F_v$。如果 $v$ 在 $F$ 中仍为素理想，则 $F_v$ 为 $k_v$ 的循环扩域；如果 $v$ 在 $F$ 中分裂，则 $F_v$ 同构于 $l$ 个 $k_v$ 的直和。在两种情形下，$F_v$ 都自然嵌入 $F_{\A,0}$。群 $G_{F_{\A,0}}$ 为限制直积 $\prod_v'G_{F_v}$。对每个 $g\in G_{F_{\A,0}}$，用 $g_v\in G_{F_v}$ 表示其 $v$ 分量。

以下固定 $\mathfrak g$ 的一个生成元 $\sigma$。设 $\boldsymbol\varphi$ 是 $C_0^\infty(G_{F_{\A,0}})$ 中具有形式 $\boldsymbol\varphi(g)=\prod_v\boldsymbol\varphi_v(g_v)$ 的元素，其中 $\boldsymbol\varphi_v\in C_0^\infty(G_{F_v})$，且除有限多个 $v$ 外，$\boldsymbol\varphi_v$ 均为 $G_{\mathfrak o(F_v)}$ 的示性函数。对每个 $v\in\mathfrak p(k)$，取 $\varphi_v\in C_0^\infty(G_{k_v})$，使 $\boldsymbol f=\boldsymbol\varphi_v$ 与 $f=\varphi_v$ 满足命题 2 的 (1)、(2)。可以假设，除有限多个 $v$ 外，$\varphi_v$ 均为 $G_{\mathfrak o_v}$ 的示性函数（见 1.3 的注记 3）。用 $\varphi$ 表示由 $\varphi(g)=\prod_v\varphi_v(g_v)$ 给出的 $G_{k_{\A,0}}$ 上的函数。由类域论（class field theory），$k^\times N_k^FF_\A^\times$ 是 $k_\A^\times$ 的指数为 $l$ 的子群。令 $\chi_0(=1),\chi_1,\ldots,\chi_{l-1}$ 为群 $k_\A^\times/(k^\times N_k^FF_\A^\times)$ 的 $l$ 个不同特征。

现在可以给出 [10] 定理 5.6 的一个阿代尔版本。

\begin{originaltheorem}{定理 2（Theorem 2）}
沿用上述记号与假设（特别是 $\boldsymbol\varphi_v$ 与 $\varphi_v$ 由命题 1 的 (1)、(2) 联系），若 $r_1,\ldots,r_n>2$，则
\begin{equation}
 l\,\operatorname{Trace}\bigl(J_\sigma T(\boldsymbol\varphi)\bigm|S(\boldsymbol\chi,r,F_\A)\bigr)
 =\sum_{i=0}^{l-1}\operatorname{Trace}\bigl(T(\varphi)\bigm|S(\chi\chi_i,r,k_\A)\bigr).
 \tag{1.10}
\end{equation}
\end{originaltheorem}
\noindent\emph{译注：所供扫描本定理 2 的括号引用为命题 1 的 (1)、(2)，而相伴函数条件在本文命题 2 中列出；此处保留扫描引用。}

\paragraph*{注记 5.}
在 [10, 定理 5.6] 中，$k$ 为有理数域，$F$ 为 $l$ 次温分歧（tamely ramified）全实阿贝尔域，并且假设 $\boldsymbol\varphi$ 非分歧。不过，该文明确描述了每个 $v$ 处的对应 $\boldsymbol\varphi_v\mapsto\varphi_v$。在这一点上，[10] 的定理 5.6 比定理 2 更精确。

\paragraph*{注记 6.}
若 $l\ne2$，则 $S(\chi\chi_i,r,k_\A)$ 与 $S(\chi,r,k_\A)$ 作为 $G_{k_{\A,0}}$-模同构。\footnote{\textbf{译者注（本次校订）：}此处应理解为经过特征扭曲后的对应，而非原有作用下的模同构。两空间的中心特征分别为 $\chi\chi_i$ 与 $\chi$，通常不同。因 $l$ 为奇数，可取 $\eta=\chi_i^{(l+1)/2}$，使 $\eta^2=\chi_i$；映射 $f(g)\mapsto\eta(\det g)f(g)$ 给出两空间的线性同构，并将右平移表示扭曲为 $\eta\circ\det$。}

\origsubsection{1.5}{}
考虑定理 2 中等式 (1.10) 两端的表示论含义。群 $G_{F_{\A,0}}$ 通过右正则表示（right regular representation）$g\mapsto T_g$ 作用于 $S(\boldsymbol\chi,r,F_\A)$。由 Jacquet--Langlands 的定理（见 [5] 的命题 10.9 和命题 11.1.1），空间 $S(\boldsymbol\chi,r,F_\A)$ 分解为不可约且互不等价的 $G_{F_{\A,0}}$-子模的代数直和（algebraic direct sum）。用 $\widehat G_{F_{\A,0}}(r,\boldsymbol\chi)$ 表示在这些不可约子模上实现的不可约表示的等价类集合。对每个 $\boldsymbol\pi\in\widehat G_{F_{\A,0}}(r,\boldsymbol\chi)$，用 $V(\boldsymbol\pi,r,\boldsymbol\chi)$ 表示实现 $\boldsymbol\pi$ 的不可约子模。于是
\[
 S(\boldsymbol\chi,r,F_\A)=\bigoplus_{\boldsymbol\pi}
 V(\boldsymbol\pi,r,\boldsymbol\chi),
\]
其中求和遍及 $\widehat G_{F_{\A,0}}(r,\boldsymbol\chi)$。记 $\boldsymbol\pi^\sigma(g)=\boldsymbol\pi(g^\sigma)$。由显然的关系 $J_\sigma T_{g^\sigma}=T_gJ_\sigma$ 可知，对每个 $\boldsymbol\pi\in\widehat G_{F_{\A,0}}(r,\boldsymbol\chi)$，都有 $\boldsymbol\pi^\sigma\in\widehat G_{F_{\A,0}}(r,\boldsymbol\chi)$，且
\[
 V(\boldsymbol\pi^\sigma,r,\boldsymbol\chi)
 =J_\sigma^{-1}V(\boldsymbol\pi,r,\boldsymbol\chi).
\]
取一个 $\boldsymbol\pi\in\widehat G_{F_{\A,0}}(r,\boldsymbol\chi)$。已知对每个 $v\in\mathfrak p(k)$，存在 $\boldsymbol\pi_v\in\widehat G_{F_v}$，使 $\boldsymbol\pi$ 等价于限制张量积（restricted tensor product）$\bigotimes_{v\in\mathfrak p(k)}'\boldsymbol\pi_v$；除有限多个 $v$ 外，$\boldsymbol\pi_v$ 都是 $G_{F_v}$ 的非分歧表示（unramified representation）。用 $\widehat G_{F_{\A,0}}(r,\boldsymbol\chi,\mathfrak g)$ 表示满足 $\boldsymbol\pi^\sigma\simeq\boldsymbol\pi$ 的元素组成的子集，其中 $\mathfrak g=\Gal(F/k)$。如果 $\boldsymbol\pi^\sigma\not\simeq\boldsymbol\pi$，则 $J_\sigma$ 在 $V(\boldsymbol\pi^\tau,r,\boldsymbol\chi)$（$\tau\in\mathfrak g$）之间诱导循环置换。因此，$J_\sigma T(\boldsymbol\varphi)$ 限制到由
\[
 \{V(\boldsymbol\pi,r,\boldsymbol\chi):
 \boldsymbol\pi\in\widehat G_{F_{\A,0}}(r,\boldsymbol\chi)
 \setminus\widehat G_{F_{\A,0}}(r,\boldsymbol\chi,\mathfrak g)\}
\]
张成的子空间上的迹为零。如果 $\boldsymbol\pi^\sigma\simeq\boldsymbol\pi$，则对每个 $v\in\mathfrak p(k)$，$\boldsymbol\pi_v^\sigma\simeq\boldsymbol\pi_v$。在 $\boldsymbol\pi_v$ 的表示空间上存在一个 $l$ 阶线性算子 $J_\sigma(\boldsymbol\pi_v)$，使
\[
 \boldsymbol\pi_v(g^\sigma)=J_\sigma(\boldsymbol\pi_v)^{-1}
 \boldsymbol\pi_v(g)J_\sigma(\boldsymbol\pi_v),\qquad g\in G_{F_v}.
\]
这样的算子只相差一个 $l$ 次单位根倍数。当 $\boldsymbol\pi_v$ 非分歧时，将 $J_\sigma(\boldsymbol\pi_v)$ 正规化，使它固定表示空间中的 $\boldsymbol\pi_v(G_{\mathfrak o_v})$-不变向量。\footnote{\textbf{译者注（本次校订）：}这里的紧子群应为 $G_{\mathfrak o(F_v)}$。$\boldsymbol\pi_v$ 是 $G_{F_v}$ 的表示，非分歧条件对应的是这一极大紧子群；$G_{\mathfrak o_v}$ 则属于基域 $k_v$。}可以假定线性算子族 $\{J_\sigma(\boldsymbol\pi_v):v\in\mathfrak p(k)\}$ 已正规化为
\[
 J_\sigma\bigm|V(\boldsymbol\pi,r,\boldsymbol\chi)
 \simeq\bigotimes_{v\in\mathfrak p(k)}J_\sigma(\boldsymbol\pi_v).
\]
对 $\boldsymbol\varphi=\prod_v\boldsymbol\varphi_v\in C_0^\infty(G_{F_{\A,0}})$，令
\[
 \boldsymbol\pi_v(\boldsymbol\varphi_v)
 =\int_{G_{F_v}}\boldsymbol\varphi_v(x)\boldsymbol\pi_v(x)\,dx.
\]
这是表示空间上的有限秩线性算子，并且
\[
 \operatorname{Trace}\bigl(J_\sigma T(\boldsymbol\varphi)
 \bigm|V(\boldsymbol\pi,r,\boldsymbol\chi)\bigr)
 =\prod_v\operatorname{trace}\bigl(J_\sigma(\boldsymbol\pi_v)
 \boldsymbol\pi_v(\boldsymbol\varphi_v)\bigr).
\]
因此
\begin{equation}
 \operatorname{Trace}\bigl(J_\sigma T(\boldsymbol\varphi)
 \bigm|S(\boldsymbol\chi,r,F_\A)\bigr)
 =\sum_{\boldsymbol\pi}\prod_{v\in\mathfrak p(k)}
 \operatorname{trace}\bigl(J_\sigma(\boldsymbol\pi_v)
 \boldsymbol\pi_v(\boldsymbol\varphi_v)\bigr),
 \tag{1.11}
\end{equation}
其中求和遍及 $\widehat G_{F_{\A,0}}(r,\boldsymbol\chi,\mathfrak g)$。类似地，对于每个 $\varphi=\prod_v\varphi_v\in C_0^\infty(G_{k_{\A,0}})$，有
\begin{equation}
 \operatorname{Trace}\bigl(T(\varphi)\bigm|S(\chi,r,k_\A)\bigr)
 =\sum_\pi\prod_{v\in\mathfrak p(k)}\operatorname{trace}\pi_v(\varphi_v),
 \tag{1.12}
\end{equation}
其中 $\pi=\bigotimes_v\pi_v$ 遍及 $\widehat G_{k_{\A,0}}(r,\chi)$，并且
\[
 \pi_v(\varphi_v)=\int_{G_{k_v}}\varphi_v(x)\pi_v(x)\,dx.
\]

回忆在等式 (1.10)、(1.11)、(1.12) 中，对所有 $v\in\mathfrak p(k)$，$\boldsymbol\varphi_v$ 和 $\varphi_v$ 均由命题 2 的等式 (1)、(2) 联系。因此，人们自然会期待：这些等式在局部给出定理 1，在整体则给出下面的定理，即 [10] 定理 3 的表示论表述。

\begin{originaltheorem}{定理 3（Theorem 3）}
假设 $r_1,\ldots,r_n>2$。

(1) 对每个 $\pi=\bigotimes_v\pi_v\in\widehat G_{k_{\A,0}}(r,\chi)$（$\pi_v\in\widehat G_{k_v}$），唯一存在 $\boldsymbol\pi=\bigotimes_v\boldsymbol\pi_v\in\widehat G_{F_{\A,0}}(r,\boldsymbol\chi,\mathfrak g)$（$\boldsymbol\pi_v\in\widehat G_{F_v}$），使对 $k$ 的每个奇素位（odd place）$v$，$\boldsymbol\pi_v$ 是 $\pi_v$ 从 $\widehat G_{k_v}$ 到 $\widehat G_{F_v}^{\mathfrak g}$ 的提升。称 $\boldsymbol\pi$ 为 $\pi$ 从 $\widehat G_{k_{\A,0}}(r,\chi)$ 到 $\widehat G_{F_{\A,0}}(r,\boldsymbol\chi,\mathfrak g)$ 的提升。

(2) 若 $l\ne2$，则对每个 $\boldsymbol\pi\in\widehat G_{F_{\A,0}}(r,\boldsymbol\chi,\mathfrak g)$，唯一存在 $\pi\in\widehat G_{k_{\A,0}}(r,\chi)$，使 $\boldsymbol\pi$ 为 $\pi$ 的提升。

(3) 若 $l=2$，则对每个 $\boldsymbol\pi\in\widehat G_{F_{\A,0}}(r,\boldsymbol\chi,\mathfrak g)$，存在 $\pi\in\widehat G_{k_{\A,0}}(r,\chi)$ 或 $\widehat G_{k_{\A,0}}(r,\chi\chi_1)$，使 $\boldsymbol\pi$ 为 $\pi$ 的提升；这里 $\chi_1$ 是类域论中对应于 $F$ 的 $k_\A^\times/k^\times$ 的二阶特征。

此外，$\pi_1,\pi_2\in\widehat G_{k_{\A,0}}(r,\chi)$ 到 $\widehat G_{F_{\A,0}}(r,\boldsymbol\chi,\mathfrak g)$ 的提升相同，当且仅当 $\pi_1=\pi_2$ 或 $\pi_1=\pi_2\otimes\chi_1(\det)$。
\end{originaltheorem}

\origsection{2}{}
本节先证明命题 2，再说明定理 2 为何支持定理 1 和定理 3。至于定理 2，只需对 [10] 中定理 1 和定理 5.6 的证明作适当调整。

\origsubsection{2.1}{}
以下三个小节沿用 1.1 和 1.3 的记号。特别地，$k$ 是 $p$ 进域。假定 $F$ 同构于 $l$ 个 $k$ 的直和，于是可以假定 $G_F$ 为 $l$ 个 $G_k$ 的直积，且对每个 $x=(x_1,\ldots,x_l)\in G_F$，有 $x^\sigma=(x_l,x_1,\ldots,x_{l-1})$。对任意 $\boldsymbol f$ 及任意 $x\in G_k$，令
\begin{equation}
 f(x)=\int_{(G_k)^{l-1}}
 \boldsymbol f\bigl((x_2\cdots x_l)^{-1}x,x_2,\ldots,x_l\bigr)
 \,dx_2\cdots dx_l.
 \tag{2.1}
\end{equation}
容易看出 $f\in C_0^\infty(G_k)$，且对任意 $c\in\mathcal C'_\sigma(G_F)$，都有 $\Lambda_\sigma(\boldsymbol f,c)=\Lambda(f,N(c))$。因此，在这一情形中命题 2 的证明十分直接。

\origsubsection{2.2}{}
下面汇总 $G_k$ 上轨道积分（orbital integral）的若干已知结果。用 $q_k$ 表示 $k$ 的剩余域的基数，用 $\varpi$ 表示 $\mathfrak o(k)$ 的极大理想的一个生成元。用 $Q_k$ 表示 $k$ 上二维半单代数的同构类集合。对每个 $K\in Q_k$，选定一个 $K$ 到 $M(2,k)$ 的 $k$-代数嵌入，并借此将 $K$ 视为 $M(2,k)$ 的子代数。对每个 $x\in K$，用 $x'$ 表示它在 $K$ 关于 $k$ 的唯一非平凡 $k$-代数自同构下的像。众所周知
\[
 \mathcal C'(G_k)=\coprod_{K\in Q_k}\coprod_t(t)^{G_k},
\]
其中对 $t$ 的并取遍 $K^\times\setminus k^\times$ 关于 $K$ 在 $k$ 上的自同构群作用的一组完全代表元。此外，
\[
 \mathcal C(G_k)\setminus\mathcal C'(G_k)
 =\coprod_{z\in k^\times}
 \left\{\begin{pmatrix}z&0\\0&z\end{pmatrix}\right\}
 \sqcup\left(z\begin{pmatrix}1&1\\0&1\end{pmatrix}\right)^{G_k}.
\]

令 $\mathfrak o(K)$ 为 $K$ 的极大 $\mathfrak o(k)$-序（maximal $\mathfrak o(k)$-order）。对每个非负整数 $m$，记 $\mathfrak o(K)_m$ 为满足 $[\mathfrak o(K):\mathfrak o(K)_m]=q_k^m$ 的唯一 $\mathfrak o(k)$-序，用 $\mathfrak o(K)_m^\times$ 表示其单位群。对每个 $t\in\mathfrak o(K)^\times\setminus\mathfrak o(k)^\times$，唯一存在非负整数 $i$，使 $t\in\mathfrak o(K)_i^\times\setminus\mathfrak o(K)_{i+1}^\times$；令 $i(t)=i$。对每个 $f\in C_0^\infty(G_k)$，令
\[
 \alpha(z)=f(z),\qquad
 \beta(z)=\Lambda\left(f,z\begin{pmatrix}1&1\\0&1\end{pmatrix}\right),\qquad z\in k^\times
\]
（参见 (1.5)）。于是 $\alpha,\beta$ 都是 $k^\times$ 上的局部常值紧支集函数。此外，对每个 $K\in Q_k$，令 $\varphi_K$ 为 $K^\times\setminus k^\times$ 上的函数 $\varphi_K(t)=\Lambda(f,t)$；注意 $\varphi_K$ 不依赖于 $K$ 到 $M(2,k)$ 的嵌入的选择。则 $\varphi_K$ 在 $K^\times\setminus k^\times$ 上局部常值。下述引理是 [7] 引理 6.2 的一个版本（亦见 [12] 的定理 2.1.1 和定理 2.2.2）。这里给出证明。

\begin{originaltheorem}{引理 2（Lemma 2）}
沿用上述记号与假设，则三元组 $\{\varphi_K\ (K\in Q_k),\alpha,\beta\}$ 满足：

(1) $\varphi_K$ 的支集在 $K^\times$ 中相对紧。

(2) 对所有 $t\in K^\times\setminus k^\times$，$\varphi_K(t)=\varphi_K(t')$。

(3) 对每个 $z\in k^\times$，存在 $z$ 的一个邻域 $U(z)$，使 $z^{-1}U(z)\subset\mathfrak o(K)^\times$，并使对任意 $t\in U(z)\setminus(U(z)\cap k^\times)$ 均有
\begin{equation}
 \varphi_K(t)=\left\{1-\frac{C(K)}{q_k-1}\right\}\alpha(z)
 +C(K)q_k^{i(tz^{-1})-1}\beta(z),
 \tag{2.2}
\end{equation}
其中 $C(K)=[\mathfrak o(K)^\times:\mathfrak o(K)_1^\times]$。
\end{originaltheorem}

\begin{proof}
选择 $\omega\in\mathfrak o(K)$，使 $\{1,\omega\}$ 为 $\mathfrak o(K)$ 的 $\mathfrak o(k)$-基，则 $\{1,\varpi^m\omega\}$ 为 $\mathfrak o(K)_m$ 的一个 $\mathfrak o(k)$-基（$m=0,1,2,\ldots$）。对每个 $t\in K$，唯一存在
\[
 \iota(t)=\begin{pmatrix}t_{11}&t_{12}\\t_{21}&t_{22}\end{pmatrix},
 \qquad\begin{pmatrix}t\\t\omega\end{pmatrix}
 =\iota(t)\begin{pmatrix}1\\\omega\end{pmatrix}.
\]
于是 $t\mapsto\iota(t)$ 是 $K$ 到 $M(2,k)$ 的一个 $k$-代数嵌入。已知 $K$ 中每个真 $\mathfrak o(K)_m$-理想都是主理想（见 [3] 的命题 1）。因此
\[
 G_k=\coprod_{m=0}^\infty G_{\mathfrak o(k)}
 \begin{pmatrix}1&0\\0&\varpi^m\end{pmatrix}\iota(K^\times).
\]
于是对任意 $f\in C_0^\infty(G_k)$ 和 $t\in K^\times\setminus k^\times$，有
\begin{equation}
 \Lambda(f,t)=\sum_{m=0}^\infty c_m\widetilde f(\iota(t)_m),
 \tag{2.3}
\end{equation}
其中
\[
 \widetilde f(x)=\int_{G_{\mathfrak o(k)}}f(uxu^{-1})\,du,
 \qquad c_m=[\mathfrak o(K)^\times:\mathfrak o(K)_m^\times],
\]
\[
 \iota(t)_m=\begin{pmatrix}1&0\\0&\varpi^m\end{pmatrix}
 \iota(t)\begin{pmatrix}1&0\\0&\varpi^{-m}\end{pmatrix}
\]
（参见 [5] 引理 7.3.2 的证明）。只需证明 $z=1$ 的情形。由于 $f$ 局部常值且紧支，存在 $1$ 在 $\mathfrak o(K)^\times$ 中的一个邻域 $U$，使对所有 $t\in U\setminus(U\cap k^\times)$ 有
\[
 \widetilde f(\iota(t)_m)
 =\widetilde f\begin{pmatrix}1&\varpi^{-m}t_{12}\\0&1\end{pmatrix}
 =\widetilde f\begin{pmatrix}1&\varpi^{i(t)-m}\\0&1\end{pmatrix},
 \qquad m=0,1,2,\ldots,
\]
以及
\[
 \widetilde f\begin{pmatrix}1&\varpi^{i(t)+n}\\0&1\end{pmatrix}
 =f(1_2),\qquad n=0,1,2,\ldots.
\]
注意 $\alpha(1)=f(1_2)$，
\[
 \beta(1)=\sum_{n\in\Z}q_k^{-n}
 \widetilde f\begin{pmatrix}1&\varpi^n\\0&1\end{pmatrix},
\]
并且
\[
 [\mathfrak o(K)^\times:\mathfrak o(K)_{m+1}^\times]
 =q_k[\mathfrak o(K)^\times:\mathfrak o(K)_m^\times],
 \qquad m=1,2,\ldots.
\]
故对任意 $t\in U\setminus(U\cap k^\times)$，有
\[
 \varphi_K(t)=\left\{1-\frac{C(K)}{q_k-1}\right\}\alpha(1)
 +C(K)q_k^{i(t)-1}\beta(1).
\]
\end{proof}

若由 $k^\times$ 上的局部常值紧支集函数 $\alpha,\beta$，以及 $K^\times\setminus k^\times$ 上的局部常值函数族 $\{\varphi_K:K\in Q_k\}$ 构成的三元组 $\{\alpha,\beta,\varphi_K:K\in Q_k\}$ 满足引理 2 的 (1)、(2)、(3)，则称它为容许三元组（admissible triple）。

下述引理也是 [7] 引理 6.2 的一个版本，可由 [12] 的推论 1.1.4 以及上一个引理推出。

\begin{originaltheorem}{引理 3（Lemma 3）}
设 $\{\alpha,\beta,\varphi_K:K\in Q_k\}$ 为容许三元组。则存在 $f\in C_0^\infty(G_k)$，使对任意 $t\in K^\times\setminus k^\times$，$\varphi_K(t)=\Lambda(f,t)$。此外，对这样的 $f$，有
\[
 \alpha(z)=f(z1_2),\qquad
 \beta(z)=\Lambda\left(f,z\begin{pmatrix}1&1\\0&1\end{pmatrix}\right).
\]
\end{originaltheorem}

\origsubsection{2.3}{}
设 $F$ 是 $k$ 的素数次 $l$ 的循环扩域。对每个 $K\in Q_k$，令 $L=K\otimes_kF$，则 $L$ 为 $F$ 上二维半单代数。选定的 $K$ 到 $M(2,k)$ 的嵌入自然延拓为 $L$ 到 $M(2,F)$ 的嵌入。此外，从 $F$ 到 $k$ 的范数映射 $N_k^F$ 自然延拓为从 $L$ 到 $K$ 的范数映射 $N_K^L$。令
\[
 L_1=\{t\in L^\times:N_K^Lt=1\},\qquad F_1=L_1\cap F,
 \qquad L'=\{t\in L^\times:N_K^Lt\in K^\times\setminus k^\times\}.
\]
下面关于 $\mathcal C_\sigma(G_F)$ 的描述出自 H. Saito（见 [10] 的引理 3.5 和注记 3.8）。

\begin{originaltheorem}{引理 4（Lemma 4）}
(1) 有
\[
 \mathcal C'_\sigma(G_F)=\bigcup_{K\in Q_k}\bigcup_t(t)^{G_F,\sigma},
\]
其中对 $t$ 的并取遍 $L'/L_1$ 关于 $L$ 在 $F$ 上的自同构群作用的一组完全代表元。

(2) 若 $l\ne2$，则
\[
 \mathcal C'_\sigma(G_F)\setminus\mathcal C_\sigma(G_F)
 =\bigcup_{z\in F^\times/F_1}(z1_2)^{G_F,\sigma}
 \cup\bigcup_{z\in F^\times/F_1}
 \left(z\begin{pmatrix}1&1\\0&1\end{pmatrix}\right)^{G_F,\sigma}.
\]
若 $l=2$，则
\[
 \mathcal C_\sigma(G_F)\setminus\mathcal C'_\sigma(G_F)
 =\bigcup_{z\in k^\times}
 \begin{pmatrix}0&1\\z&0\end{pmatrix}^{G_F,\sigma}
 \cup\bigcup_{z\in F^\times/F_1}
 \left(z\begin{pmatrix}1&1\\0&1\end{pmatrix}\right)^{G_F,\sigma}.
\]
\end{originaltheorem}

\noindent\emph{译注：所供扫描本引理 4(2) 在 $l\ne2$ 时，左端次序为 $\mathcal C'_\sigma(G_F)\setminus\mathcal C_\sigma(G_F)$；依正则类是全部类的子集这一约定，通常应反向作差。这里保留扫描次序。}

对每个 $\boldsymbol f\in C_0^\infty(G_F)$，如下定义 $k^\times$ 上的函数 $\alpha,\beta$。若 $z\in N_k^FF^\times$，写 $z=N_k^F\boldsymbol z$（$\boldsymbol z\in F^\times$），并令
\[
 \alpha(z)=\Lambda_\sigma(\boldsymbol f,\boldsymbol z1_2),\qquad
 \beta(z)=\Lambda_\sigma\left(\boldsymbol f,
 \boldsymbol z\begin{pmatrix}1&1\\0&1\end{pmatrix}\right)|l|_k
\]
（参见 (1.3)）。对 $z\in k^\times\setminus N_k^FF^\times$，令 $\beta(z)=0$；若 $l\ne2$，则令 $\alpha(z)=0$；若 $l=2$，则令
\[
 \alpha(z)=-(q_k-1)\Lambda_\sigma\left(\boldsymbol f,
 \begin{pmatrix}0&1\\z&0\end{pmatrix}\right).
\]
由引理 1 与引理 4 可知，$\alpha,\beta$ 是 $k^\times$ 上良定义的函数，并且两者均局部常值且紧支。

对每个 $K\in Q_k$，如下定义 $K^\times\setminus k^\times$ 上的函数 $\varphi_K$。若 $t\in(K^\times\setminus k^\times)\cap N_K^LL$，写 $t=N_K^L\boldsymbol t$（$\boldsymbol t\in L^\times$），令 $\varphi_K(t)=\Lambda_\sigma(\boldsymbol f,\boldsymbol t)$；若 $t\notin(K^\times\setminus k^\times)\cap N_K^LL^\times$，则令 $\varphi_K(t)=0$。引理 1 与引理 4 再次保证 $\varphi_K$ 在 $K^\times\setminus k^\times$ 上良定义。

这样，证明命题 2 就归结为证明下面的引理 5。

\begin{originaltheorem}{引理 5（Lemma 5）}
沿用上述记号，则 $\{\alpha,\beta,\varphi_K:K\in Q_k\}$ 为容许三元组。
\end{originaltheorem}

对 $K^\times$ 的一个子群 $V$（$K\in Q_k$），令
\begin{equation}
 W(V)=V\cup\left\{t\begin{pmatrix}1&1\\0&1\end{pmatrix}:t\in V\cap k^\times\right\}.
 \tag{2.4}
\end{equation}
下面的引理将用于证明引理 5。

\begin{originaltheorem}{引理 6（Lemma 6）}
对给定的 $K\in Q_k$ 和 $G_F$ 的紧子集 $C_1$，存在 $K^\times$ 的开子群 $V$ 及 $G_F$ 的紧子集 $C_2$，使得：若对某个 $w\in W(V)$ 有 $g^\sigma wg^{-1}\in C_1$，则 $g^\sigma g^{-1}\in C_2$。
\end{originaltheorem}

\begin{proof}[引理 6 的证明]
可选取 $K^\times$ 的紧开子群 $V_1$，使其满足：映射 $x\mapsto x^l$ 是 $V_1$ 到 $V_1^l$ 的同构，其逆映射为
\begin{equation}
 x\longmapsto x^{1/l}=\sum_{m=0}^\infty\binom{1/l}{m}(x-1)^m
 \tag{2.5}
\end{equation}
（$V_1$ 足够小，使该级数在 $V_1$ 上绝对且一致收敛）。用 $C_3$ 表示 $C_1$ 在范数映射 $x\mapsto N(x)$ 下的像（参见 1.1）。由于 $N(g^\sigma wg^{-1})=gN(w)g^{-1}$，若对某个 $w\in W(V_1)$ 有 $g^\sigma wg^{-1}\in C_1$，则 $gw^lg^{-1}\in C_3$。令
\[
 C_4=C_3\cap\{gw^lg^{-1}:w\in W(V_1),\ g\in G_F\}.
\]
二项式级数（binomial series）(2.5) 在 $C_4$ 上绝对且一致收敛。因此 $C_4$ 在映射 $x\mapsto x^{1/l}$ 下的像的闭包 $C_5$ 是紧的。容易看出，当 $w\in W(V_1)$ 时，$gwg^{-1}=(gw^lg^{-1})^{1/l}$。故若对某个 $w\in W(V_1)$ 有 $g^\sigma wg^{-1}\in C_1$，则 $g^\sigma g^{-1}\in C_1C_5^{-1}$。于是可取 $C_2=C_1C_5^{-1}$，$V=V_1$。
\end{proof}

\begin{proof}[引理 5 的证明]
容易看出 $\varphi_K$ 满足条件 (1)、(2)。对 $t_0\in N_k^FF^\times$，取 $z\in F^\times$，使 $t_0=N_k^Fz$。令 $C_0$ 为 $\boldsymbol f$ 的支集，选择 $K^\times$ 的一个紧开子群 $V$，使其对 $C_1=z^{-1}C_0$ 具有引理 6 所述性质。于是存在 $G_F$ 的紧子集 $C_2$，使对 $w\in W(V)$，若 $zg^\sigma wg^{-1}\in C_0$，则 $g^\sigma g^{-1}\in C_2$（$W(V)$ 由 (2.4) 给出）。

设 $\{y_i:i\in I\}$ 为双陪集（double coset）$G_{\mathfrak o(F)}\backslash G_F/G_k$ 的一组完全代表元。由引理 1 可知，映射 $g\mapsto g^\sigma g^{-1}$ 是从 $G_F/G_k$ 到 $G_F$ 的闭子集 $\{g\in G_F:N(g)=1\}$ 的同胚（homeomorphism）。因此存在 $I$ 的有限子集 $I_0$，使
\begin{equation}
 \left.
 zg^\sigma wg^{-1}\in C_0\text{，其中某个 }w\in W(V)
 \ \right\}\Longrightarrow
 g\in\bigcup_{i\in I_0}G_{\mathfrak o(F)}y_iG_k.
 \tag{2.6}
\end{equation}
当 $V$ 足够小时，对 $t\in V\setminus(V\cap k^\times)$，有 $Z_\sigma(t)=K^\times$，并且
\[
 \varphi_K(t_0t^l)=\Lambda_\sigma(\boldsymbol f,zt)
 =\sum_{i\in I_0}\mu_i\Lambda(f_i,t),
\]
其中 $\mu_i^{-1}$ 是 $G_k\cap y_i^{-1}G_{\mathfrak o(F)}y_i$ 在 $G_k$ 中的不变体积，$f_i\in C_0^\infty(G_k)$ 定义为
\[
 f_i(g)=\int_{G_{\mathfrak o(F)}}
 \boldsymbol f(zu^\sigma y_i^\sigma gy_i^{-1}u^{-1})\,du.
\]
由于 $I_0$ 有限，由引理 2 可知，$1$ 在 $K^\times$ 中存在一个邻域 $V_1\subset V$，使对每个 $t\in V_1\setminus(V_1\cap k^\times)$ 和 $i\in I_0$，均有
\[
 \Lambda(f_i,t)=\left\{1-\frac{C(K)}{q_k-1}\right\}f_i(1)
 +C(K)q_k^{i(t)-1}\Lambda\left(f_i,\begin{pmatrix}1&1\\0&1\end{pmatrix}\right).
\]
而 (2.6) 蕴含
\[
 \sum_{i\in I_0}f_i(1)\mu_i=\Lambda_\sigma(\boldsymbol f,z)=\alpha(t_0),
\]
以及
\[
 \begin{aligned}
 \sum_{i\in I_0}\mu_i\Lambda\left(f_i,\begin{pmatrix}1&1\\0&1\end{pmatrix}\right)
 &=\Lambda_\sigma\left(\boldsymbol f,z\begin{pmatrix}1&1\\0&1\end{pmatrix}\right)\\
 &=|l|_k^{-1}\beta(t_0).
 \end{aligned}
\]
还存在更小的邻域 $V_2\subset V_1$，使对 $t\in V_2\setminus(V_2\cap k^\times)$ 有 $i(t^l)=\operatorname{ord}l+i(t)$。令 $U=\{t_0t^l:t\in V_2\}$，则 $U$ 为 $t_0$ 在 $K^\times$ 中的一个邻域。我们已证明，对任意 $t\in U\setminus(U\cap k^\times)$，有
\begin{equation}
 \varphi_K(t)=\left\{1-\frac{C(K)}{q_k-1}\right\}\alpha(t_0)
 +C(K)q_k^{i(tt_0^{-1})-1}\beta(t_0).
 \tag{2.7}
\end{equation}

现在取 $t_0\in k^\times\setminus N_k^FF^\times$，令 $L=F\otimes_kK$。若 $l\ne2$，则 $N_K^LL^\times\cap k^\times=N_k^FF^\times$。因此，$t_0$ 在 $K$ 中存在一个邻域 $U$，使 $U\cap N_K^LL^\times=\varnothing$。故对每个 $t\in U\setminus(U\cap k^\times)$，都有 $\varphi_K(t)=0$。由于 $\alpha(t_0)=\beta(t_0)=0$，条件 (3) 得证。接着假设 $l=2$。若 $K$ 不是域，则 $t_0$ 在 $K^\times$ 中存在邻域 $U$，使 $U\cap N_K^LL^\times=\varnothing$，从而对每个 $t\in U\setminus(U\cap k^\times)$ 均有 $\varphi_K(t)=0$。由于 $\beta(t_0)=0$ 且 $C(K)=q_k-1$，(2.7) 对这些 $t$ 都成立。

现在假定 $K$ 是 $k$ 的二次扩域。于是存在 $s_0\in L^\times$，使 $t_0=N_K^Ls_0$。此外，$Z_\sigma(s_0)$ 同构于 $k$ 上四元数除代数（division quaternion algebra）的乘法群。

下述子引理的证明与引理 6 的证明十分相似。

\begin{originaltheorem}{子引理（Sublemma）}
沿用上述记号与假设，则对给定的 $G_F$ 的紧子集 $C_1$，存在 $K^\times$ 的开子群 $V$ 及 $G_F$ 的紧子集 $C_2$，使得：若对某个 $v\in V$ 有 $g^\sigma s_0vg^{-1}\in C_1$，则 $g^\sigma s_0g^{-1}\in C_2$。
\end{originaltheorem}

将 $Z_\sigma(s_0)$ 与 $K^\times$ 的不变测度正规化，使各自的极大紧子群的体积均等于 $1$。

由子引理容易推出，$s_0$ 在 $L^\times$ 中存在一个邻域 $U_1$，使以下积分对 $s\in U_1$ 绝对收敛，并在 $U_1$ 上给出局部常值函数：
\[
 \int_{G_F/K^\times}\boldsymbol f(g^\sigma sg^{-1})\,d\dot g.
\]
若 $t=N_K^Ls\notin k^\times$，则该积分等于 $\varphi_K(t)=\Lambda_\sigma(\boldsymbol f,s)$。若 $s=s_0$，则此积分等于
\[
 \Lambda_\sigma(\boldsymbol f,s_0)
 \int_{Z_\sigma(s_0)/K^\times}d\dot x,
\]
其中 $d\dot x$ 是由 $Z_\sigma(s_0)$ 的不变测度除以 $K^\times$ 的不变测度得到的商测度（quotient measure）。当 $K$ 分歧或非分歧时，$Z_\sigma(s_0)/K^\times$ 的体积分别等于 $1$ 或 $2$。另一方面，$C(K)$ 分别等于 $q_k$ 或 $q_k+1$。由于 $\alpha(t_0)=-(q_k-1)\Lambda_\sigma(\boldsymbol f,s_0)$ 且 $\beta(t_0)=0$，由此证明 $t_0$ 存在一个邻域 $U$，使 (2.7) 对任意 $t\in U\setminus(U\cap k^\times)$ 成立。引理 5 的证明完成。
\end{proof}

\paragraph*{注记.}
引理 5 的证明还给出：当 $F$ 为 $k$ 的扩域时，命题 2 中的 $\boldsymbol f$ 和 $f$ 满足下述关系。如果 $z=N_k^F\boldsymbol z$（$\boldsymbol z\in F^\times$），则
\[
 f(z1_2)=\Lambda_\sigma(\boldsymbol f,\boldsymbol z1_2),
\]
并且
\[
 \Lambda\left(f,z\begin{pmatrix}1&1\\0&1\end{pmatrix}\right)
 =|l|_k\Lambda_\sigma\left(\boldsymbol f,
 \boldsymbol z\begin{pmatrix}1&1\\0&1\end{pmatrix}\right).
\]
若 $z\in k^\times$ 而 $z\notin N_k^FF^\times$，则
\[
 \Lambda\left(f,z\begin{pmatrix}1&1\\0&1\end{pmatrix}\right)=0,
\]
并且
\[
 f(z1_2)=\begin{cases}
 0,&l\ne2,\\
 -(q_k-1)\Lambda_\sigma\left(\boldsymbol f,\begin{pmatrix}0&1\\z&0\end{pmatrix}\right),&l=2.
 \end{cases}
\]

\origsubsection{2.4}{}
当 $F$ 同构于 $l$ 个 $k$ 的直和时，(2.1) 明确给出了命题 2 中的对应 $\boldsymbol f\mapsto f$。当 $F$ 为 $k$ 的 $l$ 次非分歧循环扩域时，如果 $\boldsymbol f$ 关于 $G_{\mathfrak o(F)}$ 双不变，则也可以明确给出此对应。

确切地说，用 $L_0(G_k,G_{\mathfrak o(k)})$ 表示所有关于 $G_{\mathfrak o(k)}$ 左右不变的函数 $f\in C_0^\infty(G_k)$ 的集合。它关于卷积乘法（convolution product）成为交换代数。对不定元 $X,Y$，用 $\C[X,Y,X^{-1},Y^{-1}]_0$ 表示 $\C[X,Y,X^{-1},Y^{-1}]$ 中关于 $X,Y$ 对称的元素所成的子代数。对每个 $f\in L_0(G_k,G_{\mathfrak o(k)})$，令
\[
 F(f,k)[X,Y]=\sum_{m,n\in\Z}C_{mn}X^mY^n,
\]
其中
\[
 C_{mn}=q_k^{-(m-n)/2}\int_k
 f\left(\begin{pmatrix}\varpi^m&0\\0&\varpi^n\end{pmatrix}
 \begin{pmatrix}1&x\\0&1\end{pmatrix}\right)\,dx.
\]
已知映射 $f\mapsto F(f,k)$ 是从 $L_0(G_k,G_{\mathfrak o(k)})$ 到 $\C[X,Y,X^{-1},Y^{-1}]_0$ 的代数同构（参见 [11] 的定理 3）。

设 $F$ 为 $k$ 的 $l$ 次非分歧循环扩域。对每个 $\boldsymbol f\in L_0(G_F,G_{\mathfrak o(F)})$，唯一存在 $\lambda(\boldsymbol f)\in L_0(G_k,G_{\mathfrak o(k)})$，使
\[
 F(\boldsymbol f,F)(X^l,Y^l)=F(\lambda(\boldsymbol f),k)(X,Y).
\]
映射 $\boldsymbol f\mapsto\lambda(\boldsymbol f)$ 是从 $L_0(G_F,G_{\mathfrak o(F)})$ 到 $L_0(G_k,G_{\mathfrak o(k)})$ 的 $\C$-代数同态。

[6] 在更一般的背景下讨论了同态 $\lambda$。下述命题隐含于 Saito [10] 中；事实上，[10] 中 3.4--3.13 和 5.1--5.4 的相当一部分内容都是用来证明命题 3 的。

\begin{originaltheorem}{命题 3（Proposition 3）}
沿用上述记号与假设，则 $\boldsymbol f\in L_0(G_F,G_{\mathfrak o(F)})$ 与 $\lambda(\boldsymbol f)\in L_0(G_k,G_{\mathfrak o(k)})$ 由命题 2 的 (1)、(2) 联系。
\end{originaltheorem}

\origsubsection{2.5}{}
令 $M$ 为 $\mathfrak p(k)$ 的有限子集，包含 $k$ 的所有在 $F$ 中分歧的素理想。用 $C_0^\infty(G_{F_{\A,0}},M)$ 表示 $C_0^\infty(G_{F_{\A,0}})$ 中，由满足对每个 $v\notin M$ 均有 $\boldsymbol\varphi_v\in L_0(G_{F_v},G_{\mathfrak o(F_v)})$ 的 $\boldsymbol\varphi=\prod_v\boldsymbol\varphi_v$（$\boldsymbol\varphi_v\in C_0^\infty(G_{F_v})$）张成的子空间。对 $\boldsymbol\varphi=\prod_v\boldsymbol\varphi_v\in C_0^\infty(G_{F_{\A,0}},M)$，如下选择 $\varphi=\prod_v\varphi_v\in C_0^\infty(G_{k_{\A,0}})$：

当 $v\in M$ 时，选择 $\varphi_v\in C_0^\infty(G_{k_v})$，使 $\boldsymbol\varphi_v$ 和 $\varphi_v$ 由命题 2 的 (1)、(2) 联系。

当 $v\notin M$ 时，如果 $v$ 在 $F$ 中仍为素理想，令 $\varphi_v=\lambda(\boldsymbol\varphi_v)$，其中 $\lambda$ 为 2.4 引入的从 $L_0(G_{F_v},G_{\mathfrak o(F_v)})$ 到 $L_0(G_{k_v},G_{\mathfrak o(k_v)})$ 的同态。如果 $v$ 在 $F$ 中分裂，则令 $\varphi_v$ 为将 (2.1) 右端的 $\boldsymbol f$ 换成 $\boldsymbol\varphi_v$ 所给出的 $G_{k_v}$ 上的函数。

在两种情形中，$\varphi_v\in L_0(G_{k_v},G_{\mathfrak o(k_v)})$（$v\notin M$）。对任意非分歧的 $\pi_v\in\widehat G_{k_v}$，用 $\widetilde\pi_v$ 表示其从 $G_{k_v}$ 到 $G_{F_v}$ 的提升（见命题 1）。由命题 1 的第一部分可知
\begin{equation}
 \operatorname{trace}\bigl(J_\sigma(\widetilde\pi_v)
 \widetilde\pi_v(\boldsymbol\varphi_v)\bigr)
 =\operatorname{trace}\pi_v(\varphi_v),\qquad v\notin M.
 \tag{2.8}
\end{equation}
用 $\widehat G_{F_{\A,0}}(r,\boldsymbol\chi,M,\mathfrak g)$（分别为 $\widehat G_{k_{\A,0}}(r,\chi,M)$）表示 $\widehat G_{F_{\A,0}}(r,\boldsymbol\chi,\mathfrak g)$（分别为 $\widehat G_{k_{\A,0}}(r,\chi)$）中所有满足下述条件的 $\boldsymbol\pi=\bigotimes_v\boldsymbol\pi_v$（分别为 $\pi=\bigotimes_v\pi_v$）所成的子集：对任意 $v\notin M$，$\boldsymbol\pi_v$（分别为 $\pi_v$）都是 $G_{F_v}$（分别为 $G_{k_v}$）的非分歧不可约表示。对 $\boldsymbol\varphi=\prod_v\boldsymbol\varphi_v\in C_0^\infty(G_{F_{\A,0}},M)$，由 (1.10)、(1.11)、(1.12) 得
\begin{equation}
 l\sum_{\boldsymbol\pi}\prod_v
 \operatorname{trace}\bigl(J_\sigma(\boldsymbol\pi_v)
 \boldsymbol\pi_v(\boldsymbol\varphi_v)\bigr)
 =\sum_{j=0}^{l-1}\sum_\pi\prod_v\operatorname{trace}\pi_v(\varphi_v),
 \tag{2.9}
\end{equation}
其中对 $\boldsymbol\pi$ 的求和遍及 $\widehat G_{F_{\A,0}}(r,\boldsymbol\chi,M,\mathfrak g)$，对 $\pi$ 的求和遍及 $\widehat G_{k_{\A,0}}(r,\chi\chi_j,M)$。

对每个 $\boldsymbol\pi=\bigotimes_v\boldsymbol\pi_v\in\widehat G_{F_{\A,0}}(r,\boldsymbol\chi,M,\mathfrak g)$，用 $X(\boldsymbol\pi)$ 表示 $\bigcup_{j=0}^{l-1}\widehat G_{k_{\A,0}}(r,\chi\chi_j,M)$ 中所有满足以下条件的 $\pi=\bigotimes_v\pi_v$ 所成的子集：对每个 $v\notin M$，$\pi_v$ 从 $\widehat G_{k_v}$ 到 $\widehat G_{F_v}^{\mathfrak g}$ 的提升为 $\boldsymbol\pi_v$。

结合等式 (2.8)、(2.9) 与“强重数一定理”（strong multiplicity one theorem）（见 [9] 的定理 B 与 [1] 的定理 2）给出
\begin{equation}
 l\prod_{v\in M}\operatorname{trace}\bigl(J_\sigma(\boldsymbol\pi_v)
 \boldsymbol\pi_v(\boldsymbol\varphi_v)\bigr)
 =\sum_{\pi\in X(\boldsymbol\pi)}\prod_{v\in M}
 \operatorname{trace}\pi_v(\varphi_v).
 \tag{2.10}
\end{equation}
回顾 Ann Arbor 的报告，我当时低估了这一步的难度，以为定理 1 和定理 3 都能直接由 (2.10) 推出。实际上，从这一等式到两个定理，还需要相当深入的论证。[8] 则已经证明了远比定理 1、定理 3 更广泛的结果。

\begin{thebibliography}{15}
\bibitem[1]{c25-1} W. Casselmann, \emph{On some results of Atkin and Lehner}, Math. Ann. 201 (1973), 301--314.
\bibitem[2]{c25-2} K. Doi and H. Naganuma, \emph{On the functional equation of certain Dirichlet series}, Invent. Math. 9 (1969), 1--14.
\bibitem[3]{c25-3} Y. Ihara, \emph{Hecke polynomial as congruence $\zeta$-functions in elliptic modular case}, Ann. of Math. (2) 85 (1967), 267--295.
\bibitem[4]{c25-4} H. Jacquet, \emph{Automorphic forms on GL(2). II}, Lecture Notes in Math., vol. 278, Springer-Verlag, 1972.
\bibitem[5]{c25-5} H. Jacquet and R. Langlands, \emph{Automorphic forms on GL(2)}, Lecture Notes in Math., vol. 114, Springer-Verlag, 1970.
\bibitem[6]{c25-6} R. Langlands, \emph{Problems in the theory of automorphic forms}, Lecture Notes in Math., vol. 170, Springer-Verlag, 1970, pp. 18--61.
\bibitem[7]{c25-7} R. Langlands, \emph{Modular forms and $l$-adic representations}, Lecture Notes in Math., vol. 349, Springer-Verlag, 1973, pp. 361--500.
\bibitem[8]{c25-8} R. Langlands, \emph{Base change for GL(2), the theory of Saito-Shintani with applications}, Institute for Advanced Study, Princeton, N.J., 1975.
\bibitem[9]{c25-9} T. Miyake, \emph{On automorphic forms on $\GL_2$ and Hecke operators}, Ann. of Math. (2) 94 (1971), 174--189.
\bibitem[10]{c25-10} H. Saito, \emph{Automorphic forms and algebraic extensions of number fields}, Lectures in Mathematics, vol. 8, Kyoto Univ., Kyoto, Japan, 1975.
\bibitem[11]{c25-11} I. Satake, \emph{Theory of spherical functions on reductive algebraic groups over $p$-adic fields}, Publ. Math. No. 18, Inst. Hautes Études Sci., 1963.
\bibitem[12]{c25-12} J. Shalika, \emph{A theorem on semi-simple $p$-adic groups}, Ann. of Math. (2) 95 (1972), 226--242.
\bibitem[13]{c25-13} T. Shintani, \emph{On liftings of holomorphic automorphic forms (a representation theoretic interpretation of the recent work of H. Saito)}, 美日数论研讨会非正式会议录（可向东京大学获取），Ann Arbor, 1975.
\bibitem[14]{c25-14} T. Shintani, \emph{Two remarks on irreducible characters of finite general linear groups}, J. Math. Soc. Japan 28 (1976), 396--414.
\bibitem[15]{c25-15} T. Shintani, \emph{On irreducible unitary characters of a certain group extension of GL(2,C)}, J. Math. Soc. Japan 29 (1977), 165--188.
\end{thebibliography}
\begin{flushleft}\small University of Tokyo\end{flushleft}
